\chapter{Bronchospasm}

\section*{Definition}
Bronchospasm is narrowing of the lower airways from contraction of airway muscle, swelling, or mucus. It can cause wheezing, shortness of breath, chest tightness, and cough. It is a clinical finding rather than a diagnosis by itself; spirometry can help identify variable airflow limitation when asthma is suspected.

\section*{When to See a Doctor}

\begin{itemize}
 \item Severe shortness of breath at rest, inability to speak full sentences, or use of accessory respiratory muscles
 \item Oxygen saturation less than 90\% on room air or cyanosis
 \item Peak expiratory flow rate less than 50\% of personal best or predicted
 \item Symptoms are severe, worsening, or not improving with the reliever and action plan prescribed for you
 \item History of prior intubation or ICU admission for asthma flare-up
\end{itemize}

\section*{How It Develops}
Bronchospasm can result from airway inflammation, allergic mediator release, irritants, infection, exercise, or medicines. Airway muscle contraction, swelling, and mucus can occur together and narrow airflow.

\section*{Causes}
\begin{itemize}
 \item \textbf{Asthma:} Allergic (extrinsic) or non-allergic (intrinsic) triggers including allergens, exercise, cold air, and emotional stress
 \item \textbf{Respiratory infections:} Viral bronchiolitis (RSV in children), acute bronchitis, pneumonia
 \item \textbf{Environmental:} Tobacco smoke, air pollution, occupational exposures (isocyanates, grain dust, wood dust)
 \item \textbf{Drug-induced:} Beta-blockers (including ophthalmic), NSAIDs in aspirin-exacerbated respiratory disease, cholinergic agents
 \item \textbf{Other:} GERD (vagal reflex), anaphylaxis, exercise-induced bronchoconstriction, vocal cord dysfunction
\end{itemize}

\section*{Related Symptoms}
\begin{itemize}
 \item Wheezing, shortness of breath, chest tightness, cough, tachypnea, prolonged expiratory phase, hypoxia
\end{itemize}
\textbf{Important:} Severe breathlessness with little or no wheezing (a ``silent chest'') can indicate critically reduced airflow and impending respiratory failure; seek emergency care.

\section*{Medical Evaluation}
\begin{itemize}
 \item Spirometry with bronchodilator testing to look for variable airflow limitation, interpreted using age- and test-appropriate criteria
 \item Peak expiratory flow monitoring for home assessment and asthma action plan guidance
 \item Methacholine challenge test for diagnostic uncertainty when spirometry is normal
 \item Exhaled nitric oxide (FeNO) measurement as a marker of eosinophilic airway inflammation
 \item Allergy testing (skin prick or specific IgE) to identify modifiable triggers
\end{itemize}

\section*{Treatment Options}
\begin{itemize}
 \item An inhaled reliever selected as part of an individualized asthma plan; current guidelines generally prefer an inhaled corticosteroid--formoterol reliever for many adults and adolescents with asthma when available and appropriate
 \item Inhaled corticosteroid-containing treatment to reduce exacerbation risk when asthma is diagnosed
 \item Long-acting bronchodilators are used only in appropriate combinations and never as a substitute for inhaled corticosteroid treatment in asthma
 \item Leukotriene receptor antagonists (montelukast) for allergic asthma and exercise-induced bronchospasm
 \item Systemic corticosteroids for a moderate or severe asthma exacerbation when prescribed or given in urgent care
\end{itemize}

\section*{Self-Care and Home Management}
\begin{itemize}
 \item Follow your written asthma action plan for reliever use; if severe symptoms do not respond promptly, seek emergency care rather than repeatedly self-treating
 \item Identify and avoid known triggers (allergens, smoke, cold air, exercise in extreme conditions)
 \item Sit upright, lean forward on a table (tripod position) to maximize diaphragmatic excursion
 \item Pursed-lip breathing technique to maintain positive airway pressure and reduce air trapping
 \item Maintain a written asthma action plan with green/yellow/red zone parameters for daily management
\end{itemize}

\section*{References}
\begin{thebibliography}{99}
\bibitem{bronchospasm:ref1} Global Initiative for Asthma. ``Global Strategy for Asthma Management and Prevention.'' GINA, 2024.
\bibitem{bronchospasm:ref2} Papi A, Brightling C, Pedersen SE, et al. ``Asthma.'' Lancet. 2018;391(10122):783--800.
\bibitem{bronchospasm:ref3} National Asthma Education and Prevention Program. ``Expert Panel Report 3: Guidelines for the Diagnosis and Management of Asthma.'' NHLBI, 2007.
\bibitem{bronchospasm:ref4} Fanta CH. ``Asthma.'' N Engl J Med. 2009;360(10):1002--1014.
\bibitem{bronchospasm:ref5} Wechsler ME, Szefler SJ, Ortega VE, et al. ``Leukotriene Receptor Antagonists vs Inhaled Corticosteroids for Asthma.'' JAMA. 2022;327(11):1045--1055.
\bibitem{bronchospasm:ref6} O'Byrne PM, Pedersen S, Carlsson LG, et al. ``Risks of Inhaled Corticosteroids in Asthma.'' N Engl J Med. 2023;388(6):532--543.
\end{thebibliography}
